\chapter{The Helicorder}\label{ch:heli}

A \emph{helicorder} wraps a whole day of seismic data into rows, like lines of text, so that a day fits on one screen and an
event stands out as a burst on one line. The \ui{Helicorder} page shows one channel of one day; a click on any row opens a detail
view of that moment.

\begin{figure}[tbp]
\centering
\shot[\linewidth]{heli-day}
\caption{One day of the demonstration station. Each row is 30~minutes; rows alternate in colour. The local earthquake is the burst
in the 03:00 row, the noise bursts are the narrow spikes around 10:00 and 11:30, and the white gap in the 17:00 rows is a data gap.}\label{fig:heliday}
\end{figure}

\section{Opening the helicorder}
Choose \ui{Helicorder} in the sidebar. The page shows the newest day for which data exist on the laptop. The helicorder reads
the data that are \emph{already on your laptop}: for a day you have not downloaded yet, first use the \ui{Data} page
(Chapter~\ref{ch:data}); a day without data is shown as empty with a message.

\section{The controls}
The row of controls above the plot has, from left to right:
\begin{description}[leftmargin=3.2cm,style=nextline]
\item[Station, channel] the station and the channel to draw (the vertical \code{EHZ} first).
\item[\ \raisebox{-1pt}{$\langle$}\ date\ \raisebox{-1pt}{$\rangle$}] the day shown; the arrows move to the previous and the next day.
\item[counts / m/s] raw counts or ground velocity (Section~\ref{sec:units}; \ui{m/s} needs the instrument response).
\item[Band-pass] switches a band-pass filter on or off; the band is the one last set on the \ui{Data} page (1--10~Hz at first).
\item[Scale] makes the traces taller or shorter ($0.5\times$ to $8\times$); see below.
\end{description}

\section{Reading the day view}
Time runs from left to right along a row (0 to 30 minutes, labelled \emph{Minutes}); rows follow each other from 00:00
at the top. The labels at the \textbf{left} are UTC hours; the labels at the \textbf{right} are the same hours in your local time. A quiet
row is a thin line; shaking makes it fatter; a hole is blank.

\paragraph{The amplitude scale.} All rows share \emph{one} scale, so rows can be compared. At $1\times$ the ordinary signal fills about
one row and only the strongest one per cent exceeds it, so a single huge event does not flatten the rest of the day. If everything
looks too flat use \ui{2$\times$} or \ui{4$\times$}; if rows overlap too much use \ui{0.5$\times$}. Each row has its own offset removed first, so the
instrument's DC offset does not push traces out of their rows (Theory Manual, Section~9.2).

\section{Opening the detail view}
\textbf{Click anywhere on a row.} A panel opens above the day view showing about two minutes centred on the point you clicked.
(The click is accurate to about ten seconds; use the buttons below to adjust.)

\begin{figure}[tbp]
\centering
\annotshot[\linewidth]{heli-detail}{\calloutsHeli}
\caption{The detail view after clicking the 03:00 row at 12~minutes. (1) Station, channel, time range and units. (2) Fit.
(3) Box zoom. (4) Zoom out. (5) Zoom in. (6) Earlier. (7) Later. (8) Waveform, Spectra, Spectrogram. (9) Clipboard. (10) Close.}\label{fig:heli}
\end{figure}

\subsection{The toolbar}
\begin{description}[leftmargin=3.2cm,style=nextline]
\item[Zoom out, zoom in \hspace{0.2em}(4, 5)] show twice or half as much time around the same centre, between 5 seconds and one hour.
\item[Earlier, later \hspace{0.2em}(6, 7)] move the window by half of its length.
\item[Fit \hspace{0.2em}(2)] rescales the vertical axis to the data that are in view and removes any amplitude zoom. The time window
  is kept. To go back to the original two minutes, click the row again.
\item[Box zoom \hspace{0.2em}(3)] see below.
\item[Waveform / Spectra / Spectrogram \hspace{0.2em}(8)] choose how the window is drawn.
\item[Clipboard \hspace{0.2em}(9)] opens the Waveform Clipboard (Section~\ref{sec:clip}).
\item[Close \hspace{0.2em}(10)] hides the panel.
\end{description}

\section{Three views of the same window}
\paragraph{Waveform} is the amplitude against time. Here the P wave arrives at 03:12:20 and the larger S wave a few seconds later.

\paragraph{Spectra} shows how much of each frequency the window contains, on logarithmic axes: frequency in hertz horizontally,
amplitude vertically. It is good for finding a steady interfering frequency and for seeing where an earthquake has its energy
(Theory Manual, Chapter~5).

\begin{figure}[tbp]
\centering
\shot[0.92\linewidth]{heli-spectra}
\caption{\ui{Spectra} mode: amplitude spectrum of the same two minutes.}
\end{figure}

\paragraph{Spectrogram} shows the frequencies against time for the window: the sudden broadband arrival of the earthquake and its
fading coda stand out clearly (Theory Manual, Chapter~6).

\begin{figure}[tbp]
\centering
\shot[0.92\linewidth]{heli-spectrogram}
\caption{\ui{Spectrogram} mode.}
\end{figure}

\section{Zooming in with a rectangle}
The \ui{Box zoom} tool lets you drag a rectangle on the waveform and zoom into exactly that part, in time \emph{and} in amplitude.
\begin{enumerate}
\item Press the Box zoom button (3); it turns blue.
\item Press the mouse button at one corner of the part you want, drag to the opposite corner, and release.
\item The window now shows only the time range of the rectangle, reloaded at full detail, and the vertical axis shows only its amplitude range.
\end{enumerate}
Switch the button off to go back to normal mouse behaviour. The tool works in \ui{Waveform} mode.

\begin{figure}[tbp]
\centering
\shot[0.92\linewidth]{heli-box-drag}
\caption{Dragging a rectangle around the earthquake onset (the blue rectangle).}
\end{figure}

\begin{figure}[tbp]
\centering
\shot[0.92\linewidth]{heli-box-result}
\caption{After releasing: 29 seconds from 03:12:13 to 03:12:42. The P wave arrives at 03:12:20 and the stronger S wave at about
03:12:29. The vertical axis now shows only the amplitudes of the rectangle.}
\end{figure}

A rectangle narrower than five seconds is widened to five seconds around its centre. \ui{Fit} removes the amplitude zoom;
the zoom buttons change the time span as before.

\section{The Waveform Clipboard}\label{sec:clip}
The clipboard button (9) opens a dialog with the three views stacked: waveform, spectrogram and spectrum, for the window that is
currently on screen. The camera button at the top right saves them as \textbf{one picture} (PNG) that your browser downloads;
the picture has a title with the station, channel, time range in UTC, units and filter. The cross closes the dialog.

\begin{figure}[tbp]
\centering
\shot[0.92\linewidth]{heli-clipboard}
\caption{The Waveform Clipboard for a 20-second window. The camera saves the three panels as one picture.}
\end{figure}

\section{Tips}
\begin{itemize}
\item To find events quickly, scan the day view for fat or tall spots on a row, then click them.
\item Turn the band-pass on (1--10~Hz) to see local earthquakes more clearly; turn it off to see the real signal.
\item A very short event may be one burst on a single row; use \ui{Box zoom} on its onset to read the P and S arrival times to a
  fraction of a second.
\item Man-made noise (footsteps, a door, a vehicle) looks like a short burst without the slowly fading tail that an earthquake has.
\end{itemize}
